%======================================================================
\chapter{Introduction}\label{ch:intro}
%======================================================================

\section{What this volume is about}

The classification of homogeneous plane continua is a theorem about the \emph{plane}, and its proof
uses a good deal of plane topology: which compact sets separate the plane, how separation behaves
under unions and homeomorphisms, polygons and their insides, and covers of plane continua by
tree chains. This volume develops the part of this material that can be done with two elementary
tools, continuous logarithms and polygons. Moore's topological characterization of the plane,
decompositions of the plane and the Jordan curve theorem for arbitrary simple closed curves are
left to Volume~5.

The volume has five chapters.

\begin{itemize}
\item \textbf{Chapter~\ref{ch:log}: Continuous logarithms.} Uniqueness of logarithms up to a
constant, existence on contractible spaces from the homotopy lifting property of the exponential
map, explicit logarithms on small balls, the unicoherence of the plane and Brouwer's theorem on the
frontiers of complementary domains.
\item \textbf{Chapter~\ref{ch:sep}: Separation by logarithms; Janiszewski and Borsuk.} A compact set
$K$ separates $x$ from $y$ exactly when $(z-x)/(z-y)$ has no logarithm on $K$ (in the two directions
needed). Consequences: Janiszewski's theorem, Borsuk's theorem (a compact set with connected
complement carries a logarithm of every nonvanishing function), and the topological invariance of
plane separation.
\item \textbf{Chapter~\ref{ch:poly}: Polygons and the polygonal Jordan curve theorem.} The crossing
parity with a half-open crossing rule, the Jordan curve theorem for simple closed polygons,
polygonal paths and arcs, and the absence of a logarithm of $z-p$ on a polygon around $p$.
\item \textbf{Chapter~\ref{ch:enc}: Graphs, mod-2 cycles and enclosing polygons.} A finite graph of
segments separating two points contains a cycle, and then a simple closed polygon, separating them.
Applied to the boundary of a union of grid squares this gives \emph{enclosing polygons}: a continuum
$M$ not separating the plane lies inside a simple closed polygon in any given neighbourhood of $M$.
\item \textbf{Chapter~\ref{ch:tree}: Tree-like continua.} Bing's theorem that a plane continuum is
tree-like if and only if none of its subcontinua separates the plane; logarithms on tree-like
continua; a weak form of the invariance of plane separation; a survey of Bing's Theorem~7.
\end{itemize}

The main results, in the form used later, are:

\begin{theorem}[Janiszewski; Theorem~\ref{thm:janiszewski}]
Let $A,B\subseteq\R^2$ be compact with $A\cap B$ connected or empty. If neither $A$ nor $B$
separates $x$ from $y$, then $A\cup B$ does not separate $x$ from $y$.
\end{theorem}

\begin{theorem}[Borsuk; Theorem~\ref{thm:borsuk}]
If $K\subseteq\R^2$ is compact and $\R^2\setminus K$ is connected, every continuous nowhere zero
function on $K$ has a continuous logarithm on $K$.
\end{theorem}

\begin{theorem}[Theorem~\ref{thm:sep-invariant}]
If two compact subsets of the plane are homeomorphic and one separates the plane, so does the other.
\end{theorem}

\begin{theorem}[Jordan curve theorem for polygons; Theorem~\ref{thm:jct-poly}]
The complement of a simple closed polygon has exactly two components, one bounded and one unbounded,
and the polygon is the frontier of each.
\end{theorem}

\begin{theorem}[Enclosing polygons; Theorem~\ref{thm:enclose}]
Let $M$ be a continuum that does not separate the plane and $G\supseteq M$ open. There is a simple
closed polygon $J\subseteq G\setminus M$ whose inside contains $M$ and lies in $G$.
\end{theorem}

\begin{theorem}[Bing; Corollary~\ref{cor:treelike-iff}]
A plane continuum is tree-like if and only if none of its subcontinua separates the plane.
\end{theorem}

\section{Where these results are used}

\begin{itemize}
\item Janiszewski's theorem for compact sets is used throughout the plane-topological parts of the
proof: in Moore's characterization of the plane (Volume~5), in Hagopian's arguments on
indecomposable homogeneous plane continua (Volume~6), and in Jones's work on decomposable
homogeneous continua (Volume~7).
\item The unicoherence of the plane and Brouwer's theorem enter Volume~5 and Jones's arguments.
\item Borsuk's theorem and the invariance of plane separation (Theorem~\ref{thm:sep-invariant}) are
used for Jones's aposyndetic decomposition (Volume~7). The weak invariance of
Section~\ref{sec:weak-invariance} is the form used in Hagopian's argument (Volume~6).
\item Polygons, the polygonal Jordan curve theorem and enclosing polygons are the basis of Bing's
Theorems~6 and~7 and are used again in Volumes~5 and~6.
\end{itemize}

\section{Where the formal proof differs from the literature}\label{sec:intro-diff}

\begin{enumerate}
\item \textbf{Logarithms instead of homology.} Separation is detected by logarithms throughout; no
homology is used. The covering property
of the exponential map, Stone--Weierstrass approximation and Sard's theorem in equal dimensions are
quoted from Mathlib (Chapters~\ref{ch:log} and~\ref{ch:sep}).
\item \textbf{A half-open crossing rule.} The crossing parity of a polygon is defined by a half-open
rule which counts correctly when a ray passes through vertices or along edges, so that the proof of
the polygonal Jordan curve theorem needs no general-position perturbation (Chapter~\ref{ch:poly}).
\item \textbf{Janiszewski's theorem for compact sets, not closed sets.} The formal development
first stated Janiszewski's theorem for closed sets, and this is false: two closed half-lines
meeting in a point form a line (Chapter~\ref{ch:sep}). The correct hypothesis is compactness.
\item \textbf{Bing's Theorem~6 by the compact Janiszewski theorem.} The formal development first
proved Janiszewski's theorem only for polygonal sets, by mod-2 cycles (Remark~\ref{rem:jan-cycles}),
and derived Bing's Theorem~6 from Theorem~7, whose links are polygon interiors. These notes prove
Theorem~6 directly from Bing's notched-square covers, with the compact form of Janiszewski's
theorem; this proof was formalized for these notes (in the file
\texttt{HPC/Plane/Theorem6Direct.lean}). The polygonal form of Janiszewski's theorem is deduced from
the compact one.
\item \textbf{Definition of tree-likeness.} The links of the tree chains are required to be bounded.
Without this requirement the formal definition would have been satisfied by every continuum
(Chapter~\ref{ch:tree}).
\item \textbf{The unstated converse.} Bing states Theorem~7 for tree-like continua and proves it for
continua with no separating subcontinuum. The implication between the two hypotheses, which he
does not state, is Theorem~\ref{thm:treelike-nonsep}.
\item \textbf{Weak invariance without Borsuk.} The step ``a homeomorphic copy of a separating
continuum separates the plane'' in Hagopian's argument is replaced by a weaker statement proved from
Bing's Theorem~6 and logarithms on tree-like continua (Section~\ref{sec:weak-invariance}).
\end{enumerate}

\section{Prerequisites and conventions}\label{sec:intro-conv}

The reader should know the basic metric topology of \vone{Chapter~1}. From Volume~3 we use the
separation of two closed sets in a compact set (\vthree{Theorem~1.9}) and, for comparison only,
circular chains (\vthree{Definition~5.13}).

\begin{itemize}
\item \textbf{The plane.} $\R^2$ carries the \emph{sup metric}
$d(z,w)=\nrm{z-w}=\max(\abs{z_1-w_1},\abs{z_2-w_2})$, as in the formal development, where the plane is
the product $\R\times\R$. So the balls $B(c,r)$ and $\bar B(c,r)$ are open and closed squares, and the
sphere $\set{z:d(z,c)=r}$ is the boundary of a square. Balls are convex.
\item \textbf{Complex numbers.} We identify $z=(z_1,z_2)$ with the complex number $z_1+iz_2$
(\lean{HPC.Plane.toCpx}) and write $\abs w$ for the euclidean modulus. Then
$\nrm w\le\abs w\le2\nrm w$ (\lean{HPC.Plane.dist_le_norm_toCpx_sub},
\lean{HPC.Plane.norm_toCpx_sub_le}). The topology is the usual one.
\item \textbf{Continua.} A \emph{continuum} is a nonempty compact connected set
(\lean{HPC.Fold.Q78.IsContinuum}). \emph{Preconnected} means connected or empty; a closed set is
preconnected if and only if it is not the union of two disjoint nonempty closed sets
(\vone{(K2)}).
\item \textbf{Components.} For $A\subseteq\R^2$ and $x\notin A$, $\comp Ax$ is the component of $x$ in
$\R^2\setminus A$ (Definition~\ref{def:separates}). For a simple closed polygon $\Gamma$, $\inn\Gamma$ is
its inside, the union of the bounded components of $\R^2\setminus\Gamma$
(Definition~\ref{def:ins}).
\item \textbf{Logarithms.} A \emph{logarithm} of $f$ on a set $K$ is a function $g$ continuous on $K$
with $e^{g}=f$ on $K$ (Definition~\ref{def:log}). $\Log$ and $\Arg$ are the principal logarithm and
argument.
\end{itemize}

\paragraph{Formal counterparts.} As in the earlier volumes, many results are followed by a line
``Formal counterpart:'' naming the corresponding declaration of the Lean development (all names
begin with \lean{HPC.}, except those of Mathlib). These lines may be skipped.
Appendix~\ref{app:formal} collects them.

\paragraph{Exercises.} Each chapter ends with exercises. None of them is used later.
